\documentclass[../DoAn.tex]{subfiles}
\begin{document}

Chương~\ref{chapter:Methodology} đã trình bày các nền tảng lý thuyết và công nghệ. Chương này mô tả cách các nền tảng đó được tổ chức thành kiến trúc hệ thống, thiết kế chi tiết cho hai chức năng cốt lõi, dữ liệu và kho nhạc thực tế, giao diện lập trình ứng dụng và giao diện người dùng, quá trình đánh giá thực nghiệm có thể kiểm chứng, thời gian phản hồi đo được và phần kiểm thử chức năng hệ thống.

\section{Thiết kế kiến trúc}
\label{section:4.1}

\subsection{Kiến trúc tổng thể}
\label{subsection:4.1.1}
Hệ thống được tổ chức theo hai pha tách biệt nhằm đáp ứng yêu cầu phản hồi nhanh: pha ngoại tuyến (offline) và pha phục vụ (serving). Pha ngoại tuyến chịu trách nhiệm trích xuất mọi đặc trưng nặng: embedding âm thanh từ mô hình MuQ, các tín hiệu Valence từ lời, nhịp độ và độ to, rồi tính sẵn tọa độ cảm xúc và các ma trận biểu diễn cho toàn bộ kho nhạc. Pha phục vụ chỉ thực hiện các phép nhân ma trận và xếp hạng trên các đặc trưng đã tính sẵn, nên truy vấn của người dùng được trả về nhanh. Kiến trúc gồm ba tầng theo thứ tự phụ thuộc một chiều (Hình~\ref{fig:arch}): tầng dữ liệu và đặc trưng ở dưới cùng, tầng động cơ gợi ý ở giữa, và tầng giao diện ứng dụng trang đơn ở trên cùng. Tầng trên chỉ gọi xuống tầng dưới qua một giao diện lập trình ứng dụng, không có phụ thuộc ngược.

Việc tách hai pha mang lại nhiều lợi ích cụ thể. Về \textit{hiệu năng}, mọi chi phí tính toán đắt (chạy mô hình học sâu để trích embedding, tính tín hiệu cảm xúc) chỉ phải trả một lần ở pha ngoại tuyến; pha phục vụ chỉ còn các phép đại số tuyến tính nên đáp ứng truy vấn trong vài mili-giây (mục~\ref{subsec:latency}). Về \textit{tính tất định}, vì nhãn cảm xúc đã được đóng băng thành bảng tra cứu và không có mô hình nào được suy luận lúc phục vụ, cùng một truy vấn luôn cho cùng kết quả --- đáp ứng yêu cầu NFR04. Về \textit{khả năng triển khai}, máy chủ phục vụ không cần GPU hay bộ nhớ lớn để chạy mô hình, chỉ cần đủ RAM giữ các ma trận đặc trưng, nên có thể triển khai trên hạ tầng khiêm tốn. Về \textit{khả năng bảo trì}, ranh giới rõ ràng giữa ba tầng cho phép thay thế một thành phần (ví dụ đổi mô hình embedding âm thanh) mà không ảnh hưởng tầng giao diện, miễn là định dạng tệp đặc trưng được giữ nguyên.

\begin{figure}[h!]
\centering
\begin{tikzpicture}[font=\small,
   tier/.style={draw, rounded corners, align=center, minimum width=11cm, minimum height=1.3cm},
   arr/.style={-{Stealth[length=2.5mm]}, very thick},
   lbl/.style={font=\scriptsize\itshape, gray}]
  \node[tier, fill=green!7] (ui) at (0,4.0) {\textbf{Tầng giao diện (SPA React)}\\\footnotesize 2 giao diện: ``đắm chìm'' 3D \& ``classic'' 2D --- dùng chung logic};
  \node[tier, fill=blue!7] (eng) at (0,2.0) {\textbf{Tầng động cơ gợi ý (FastAPI + NumPy)}\\\footnotesize gợi ý theo màu, gợi ý bài tương tự, tra cứu cảm xúc};
  \node[tier, fill=black!5] (data) at (0,0.0) {\textbf{Tầng dữ liệu \& đặc trưng (tệp .npy / .json)}\\\footnotesize embedding MuQ, embedding lời, tọa độ V-A, nhịp độ, chỉ mục cover};
  \draw[arr] (ui) -- node[lbl, right] {gọi API (REST/JSON)} (eng);
  \draw[arr] (eng) -- node[lbl, right] {nạp \& tra cứu} (data);
  \node[lbl, right=0.1cm of ui.east, text width=2.2cm] {pha phục vụ};
  \node[lbl, right=0.1cm of data.east, text width=2.2cm] {pha ngoại tuyến chuẩn bị};
\end{tikzpicture}
\caption{Kiến trúc ba tầng của Brightify với phụ thuộc một chiều từ trên xuống}
\label{fig:arch}
\end{figure}

\subsection{Pha ngoại tuyến và luồng dữ liệu}
\label{subsection:4.1.2}
Pha ngoại tuyến là một chuỗi xử lý biến kho nhạc thô thành bộ đặc trưng phục vụ. Hình~\ref{fig:offline} mô tả luồng dữ liệu: từ tệp âm thanh và lời, hệ thống trích song song embedding MuQ, các tín hiệu Valence từ lời, nhịp độ sạch và độ to; các tín hiệu này được tổ hợp thành tọa độ Valence--Arousal; cuối cùng mọi ma trận biểu diễn được đóng gói thành một bản phát hành phục vụ.

\begin{figure}[h!]
\centering
\begin{tikzpicture}[font=\small,
   io/.style={draw, rounded corners, align=center, fill=black!4, text width=2.0cm, minimum height=0.8cm},
   proc/.style={draw, align=center, fill=blue!6, text width=2.4cm, minimum height=0.8cm},
   store/.style={draw, align=center, fill=green!7, text width=2.2cm, minimum height=0.8cm},
   arr/.style={-{Stealth[length=2mm]}, thick}]
  \node[io] (audio) at (0,1.4) {Tệp âm thanh};
  \node[io] (lyr) at (0,-1.4) {Lời bài hát};
  \node[proc] (muq) at (3.0,2.1) {Embedding MuQ};
  \node[proc] (bpm) at (3.0,0.7) {Nhịp độ + độ to};
  \node[proc] (vsig) at (3.0,-1.4) {Tín hiệu Valence (4 nguồn)};
  \node[proc, text width=2.2cm] (va) at (6.4,0.0) {Tổ hợp V-A (EWE + Arousal)};
  \node[store] (feat) at (9.6,0.0) {Bản phát hành đặc trưng};
  \draw[arr] (audio) -- (muq);
  \draw[arr] (audio) -- (bpm);
  \draw[arr] (lyr) -- (vsig);
  \draw[arr] (muq) -- (va);
  \draw[arr] (bpm) -- (va);
  \draw[arr] (vsig) -- (va);
  \draw[arr] (muq.east) to[out=0,in=120] (feat.north west);
  \draw[arr] (va) -- (feat);
\end{tikzpicture}
\caption{Luồng dữ liệu pha ngoại tuyến: trích đặc trưng song song, tổ hợp tọa độ cảm xúc, đóng gói bản phát hành phục vụ}
\label{fig:offline}
\end{figure}

\subsection{Pha phục vụ và triển khai}
\label{subsection:4.1.3}
Ở pha phục vụ, động cơ gợi ý nạp toàn bộ đặc trưng vào bộ nhớ một lần lúc khởi động, sau đó mỗi truy vấn chỉ là các phép nhân ma trận. Hình~\ref{fig:deploy} mô tả sơ đồ triển khai: trình duyệt chạy ứng dụng trang đơn, gọi máy chủ FastAPI qua REST; máy chủ giữ các ma trận đặc trưng trong bộ nhớ; tệp âm thanh được phục vụ tĩnh (hoặc chuyển hướng qua mạng phân phối nội dung). Bản phát hành đặc trưng do một máy trạm chạy pha ngoại tuyến sinh ra và nạp lên máy chủ.

\begin{figure}[h!]
\centering
\begin{tikzpicture}[font=\small,
   node/.style={draw, rounded corners, align=center},
   db/.style={draw, cylinder, shape border rotate=90, aspect=0.25, align=center, fill=black!4, minimum height=1.1cm, minimum width=1.8cm},
   arr/.style={-{Stealth[length=2mm]}, thick}]
  \node[node, fill=green!7, text width=2.6cm, minimum height=1.0cm] (br) at (0,0) {\textbf{Trình duyệt}\\\footnotesize SPA React};
  \node[node, fill=blue!7, text width=3.0cm, minimum height=1.0cm] (api) at (5.0,0) {\textbf{Máy chủ FastAPI}\\\footnotesize ma trận đặc trưng trong RAM};
  \node[db] (feat) at (9.4,0) {Tệp đặc trưng\\.npy/.json};
  \node[node, text width=2.4cm, minimum height=0.9cm] (work) at (9.4,-2.2) {Máy trạm\\(pha ngoại tuyến)};
  \draw[arr] (br) -- node[above, font=\scriptsize] {REST/JSON} (api);
  \draw[arr] (api) -- node[above, font=\scriptsize] {nạp} (feat);
  \draw[arr] (work) -- node[right, font=\scriptsize] {sinh \& triển khai} (feat);
  \draw[arr, dashed] (api.south) to[out=-90,in=-90,looseness=1.1] node[below, font=\scriptsize, pos=0.5] {âm thanh tĩnh / CDN} (br.south);
\end{tikzpicture}
\caption{Sơ đồ triển khai: trình duyệt $\leftrightarrow$ máy chủ FastAPI (giữ đặc trưng trong RAM); bản phát hành đặc trưng do máy trạm ngoại tuyến sinh ra}
\label{fig:deploy}
\end{figure}

Về tổ chức mã nguồn, hệ thống được chia thành các gói theo trách nhiệm (Hình~\ref{fig:component}): gói \texttt{core} chứa động cơ gợi ý và ánh xạ màu; gói con \texttt{core.ranking} chứa các bước truy hồi, nhất quán, đa dạng và hành trình; gói \texttt{api} cung cấp các endpoint; gói \texttt{frontend} là ứng dụng trang đơn; gói \texttt{data} lưu các tệp đặc trưng; và gói \texttt{tools} chứa các script ngoại tuyến và đánh giá.

\begin{figure}[h!]
\centering
\begin{tikzpicture}[font=\small,
   pkg/.style={draw, rounded corners, align=center, fill=black!3, minimum width=2.4cm, minimum height=0.85cm},
   arr/.style={-{Stealth[length=2mm]}, thick}]
  \node[pkg, fill=green!7] (fe) at (0,2.4) {\texttt{frontend} (SPA)};
  \node[pkg, fill=blue!7] (api) at (0,1.0) {\texttt{api} (FastAPI)};
  \node[pkg, fill=blue!10, minimum width=3.0cm] (core) at (-2.0,-0.6) {\texttt{core} (engine)};
  \node[pkg, minimum width=3.0cm] (rank) at (-2.0,-2.0) {\texttt{core.ranking}};
  \node[pkg] (data) at (2.6,-0.6) {\texttt{data}};
  \node[pkg] (tools) at (2.6,-2.0) {\texttt{tools}};
  \draw[arr] (fe) -- (api);
  \draw[arr] (api) -- (core);
  \draw[arr] (core) -- (rank);
  \draw[arr] (core) -- (data);
  \draw[arr] (tools) -- (data);
\end{tikzpicture}
\caption{Sơ đồ gói: phụ thuộc một chiều giữa các thành phần mã nguồn}
\label{fig:component}
\end{figure}

\subsection{Thiết kế lớp động cơ gợi ý}
\label{subsection:4.1.5}
Trung tâm của tầng động cơ là lớp \texttt{MusicRecommender}, giữ trạng thái kho nhạc và các ma trận đặc trưng, đồng thời cung cấp hai phương thức gợi ý chính. Lớp này ủy thác các bước hậu xử lý cho các hàm trong gói \texttt{core.ranking} (Hình~\ref{fig:class}). Cách tách này giữ cho \texttt{MusicRecommender} chỉ điều phối, còn từng bước (truy hồi, nhất quán, đa dạng, hành trình) là các hàm thuần có thể kiểm thử độc lập.

\begin{figure}[h!]
\centering
\begin{tikzpicture}[font=\footnotesize,
   class/.style={draw, rectangle split, rectangle split parts=2, align=left, fill=blue!5, text width=6.0cm},
   helper/.style={draw, rounded corners, align=left, fill=black!4, text width=4.6cm},
   arr/.style={-{Stealth[length=2mm]}, thick}]
  \node[class] (mr) at (0,0) {
    \textbf{MusicRecommender}
    \nodepart{two}
    + df, song\_va, muq\_emb, lyrics\_emb\\
    + color\_mapper\\
    + recommend\_by\_colors(hex, top\_k)\\
    + recommend\_by\_song(idx, top\_k)\\
    + \_precompute\_all\_features()};
  \node[helper] (rk) at (0,-3.4) {
    \textbf{core.ranking}\\
    retrieval $\cdot$ coherence $\cdot$ diversity\\
    journey $\cdot$ artist (MMR, RRF, waypoint)};
  \node[helper, text width=3.0cm] (cm) at (5.2,0) {
    \textbf{AdvancedColorMapper}\\
    hsl\_to\_va() $\cdot$ hex\_to\_hsl()};
  \draw[arr] (mr) -- node[right, font=\scriptsize] {ủy thác} (rk);
  \draw[arr] (mr) -- node[above, font=\scriptsize] {dùng} (cm);
\end{tikzpicture}
\caption{Sơ đồ lớp rút gọn: \texttt{MusicRecommender} điều phối và ủy thác hậu xử lý cho \texttt{core.ranking}}
\label{fig:class}
\end{figure}

\subsection{Mô hình dữ liệu kho đặc trưng}
\label{subsection:4.1.6}
Dữ liệu phục vụ được tổ chức quanh thực thể trung tâm là \textit{Bài hát}, mỗi bài liên kết với các đặc trưng đã tính sẵn (Hình~\ref{fig:erd}). Quan hệ cover là quan hệ nhiều--nhiều giữa các bài (1124 bài gom thành 979 cặp). Các đặc trưng nặng (embedding âm thanh, embedding lời) được lưu thành ma trận căn theo thứ tự bài, còn tọa độ cảm xúc và nhịp độ lưu theo khóa định danh bài.

\begin{figure}[h!]
\centering
\begin{tikzpicture}[font=\footnotesize,
   ent/.style={draw, rounded corners, align=center, fill=blue!5, minimum width=2.4cm, minimum height=0.8cm},
   attr/.style={draw, align=center, fill=black!3, text width=2.4cm, minimum height=0.7cm},
   rel/.style={-{Stealth[length=2mm]}, thick}]
  \node[ent] (song) at (0,0) {\textbf{Bài hát}\\\footnotesize track\_id (PK)};
  \node[attr] (muq) at (-4.2,1.4) {Embedding MuQ\\$1024$ chiều};
  \node[attr] (lyr) at (-4.2,-1.4) {Embedding lời\\$1024$ chiều};
  \node[attr] (va) at (4.2,1.4) {Tọa độ V-A\\(valence, arousal)};
  \node[attr] (bpm) at (4.2,-1.4) {Nhịp độ (BPM)\\+ độ to};
  \node[attr] (cov) at (0,-2.2) {Cặp cover\\(track\_a, track\_b)};
  \draw[rel] (song) -- node[above, font=\scriptsize] {1--1} (muq);
  \draw[rel] (song) -- node[below, font=\scriptsize] {1--1} (lyr);
  \draw[rel] (song) -- node[above, font=\scriptsize] {1--1} (va);
  \draw[rel] (song) -- node[below, font=\scriptsize] {1--1} (bpm);
  \draw[rel] (song) -- node[right, font=\scriptsize] {$n$--$n$} (cov);
\end{tikzpicture}
\caption{Lược đồ thực thể--quan hệ rút gọn của kho đặc trưng phục vụ}
\label{fig:erd}
\end{figure}

\section{Thiết kế chi tiết hai chức năng cốt lõi}
\label{section:4.2}

\subsection{Gán tọa độ cảm xúc cho bài hát}
\label{subsection:4.2.1}
Mỗi bài hát được gán một cặp Valence và Arousal, dùng chung cho cả hai chức năng. Valence được tính bằng tổ hợp EWE (Evaluator Weighted Estimator) \cite{grimm2005ewe} của bốn tín hiệu, trong đó trọng số bằng độ tin cậy đo được của từng tín hiệu (Bảng~\ref{table:ewe}). Có thể thấy ba tín hiệu lời chiếm phần lớn trọng số, còn tín hiệu âm thanh chỉ đóng vai trò bổ trợ, phản ánh đúng việc Valence nằm chủ yếu ở nội dung lời. Hình~\ref{fig:ewe-weights} trực quan hóa các trọng số này.

\begin{table}[h!]
\centering
\caption{Các tín hiệu Valence và vai trò trong tổ hợp EWE (trọng số tỉ lệ với độ tin cậy đo được của từng tín hiệu)}
\label{table:ewe}
\begin{tabular}{p{2.6cm} p{4.6cm} p{5.4cm}}
\toprule
\textbf{Tín hiệu Valence} & \textbf{Vai trò trong tổ hợp EWE} & \textbf{Nguồn} \\
\midrule
vn\_lex (từ điển) & Lời --- trọng số lớn & NRC-VAD-VN \cite{mohammad2018nrcvad} + VnEmoLex \cite{doan2022vnemolex} \\
emobank & Lời (liên ngữ) --- trọng số lớn & XLM-R \cite{conneau2020xlmr} + EmoBank \cite{buechel2017emobank} \\
vn\_sent & Lời (ngữ cảnh VN) --- trọng số vừa & ViSoBERT \cite{nguyen2023visobert} + UIT-VSMEC \cite{ho2020vsmec} \\
muq (âm thanh) & Bổ trợ --- trọng số nhỏ nhất & MuQ \cite{zhu2025muq} \\
\bottomrule
\end{tabular}
\end{table}

\begin{figure}[h!]
\centering
\includegraphics[width=0.62\textwidth]{Hinhve/fig_ewe_weights.png}
\caption{Trọng số tổ hợp Valence theo độ tin cậy (EWE): ba tín hiệu lời chiếm phần lớn, tín hiệu âm thanh bổ trợ}
\label{fig:ewe-weights}
\end{figure}

Arousal được tính từ một tổ hợp ba đặc trưng âm thanh: đặc trưng MuQ chiếm 0.574, nhịp độ (số nhịp trên phút) chiếm 0.35 và độ to chiếm 0.076. Phần trọng số giữa MuQ và độ to được hồi quy trên tập DEAM do con người gán nhãn \cite{aljanaki2017deam}, còn trọng số nhịp độ được \textit{nâng có chủ đích} (tinh chỉnh trên DEAM) vì biểu diễn MuQ nắm bắt độ to tốt nhưng nắm bắt nhịp độ yếu. Khi kiểm chứng chéo trên DEAM, tổ hợp này đạt hệ số tương quan 0.69, cao hơn mức 0.647 của biểu diễn MERT trước đó, và tương quan giữa Arousal với nhịp độ đo được đạt 0.47. Thuật toán~\ref{alg:valabel} tóm tắt toàn bộ quy trình gán nhãn.

\begin{algorithm}[h]
\caption{Gán tọa độ cảm xúc Valence--Arousal cho mỗi bài hát}
\label{alg:valabel}
\KwIn{Lời $\ell$, embedding MuQ $\mathbf{m}$, nhịp độ $t$, độ to $\rho$ của bài hát}
\KwOut{Cặp $(v, a)$ trong $[0,1]^2$}
$v_{\text{lex}} \leftarrow$ điểm từ điển NRC-VAD-VN/VnEmoLex trên $\ell$\;
$v_{\text{sent}} \leftarrow$ đầu dò ViSoBERT($\ell$);\quad $v_{\text{emo}} \leftarrow$ đầu dò XLM-R($\ell$);\quad $v_{\text{muq}} \leftarrow$ đầu dò MuQ($\mathbf{m}$)\;
$v \leftarrow \text{EWE}(v_{\text{lex}}, v_{\text{sent}}, v_{\text{emo}}, v_{\text{muq}})$ \tcp*{trọng số = độ tin cậy, Pt.~\eqref{eq:ewe}}
$a_{\text{muq}} \leftarrow$ đầu dò Arousal MuQ($\mathbf{m}$)\;
$a \leftarrow \text{chuẩn-hóa-hạng}\big(0.574\,\text{rank}(a_{\text{muq}}) + 0.35\,\text{rank}(t) + 0.076\,\text{rank}(\rho)\big)$ \tcp*{MuQ/độ-to fit DEAM; nhịp độ chỉnh tay}
\Return chuẩn hóa $(v, a)$ về $[0,1]^2$\;
\end{algorithm}

Bộ nhãn V-A này được kiểm hai tính chất. \textit{Tái lập}: chạy lại toàn bộ quy trình từ các tín hiệu nền bằng \texttt{tools/build\_labels\_repro.py} cho ra đúng bộ nhãn đang phục vụ (tương quan Spearman $\rho = 1.00$ trên cả Valence và Arousal). \textit{Hội tụ độc lập}: đối chiếu với một bộ tham chiếu V-A sinh riêng bằng mô hình GPT-4o-mini chỉ đọc lời bài hát (độ tin cậy test--retest ICC $= 0.99$), bộ nhãn đạt tương quan $\rho = 0.63$ cho Valence và $\rho = 0.41$ cho Arousal trên toàn bộ 5.138 bài --- mức đồng thuận có ý nghĩa với một bộ đánh giá hoàn toàn độc lập, trong đó Arousal thấp hơn là hợp lý vì nó được suy từ âm thanh còn GPT chỉ đọc lời.

\subsection{Ánh xạ màu sắc sang tọa độ cảm xúc}
\label{subsection:4.2.2}
Màu được chuyển sang không gian Oklab \cite{ottosson2020oklab}, sau đó Valence của màu được suy ra bằng hồi quy ridge khớp với chuẩn ICEAS \cite{jonauskaite2020universal}, đạt hệ số tương quan 0.969 so với chuẩn này. Arousal của màu được tính theo công thức suy từ quan hệ màu--âm nhạc của Whiteford \cite{whiteford2018color} trong Phương trình~\eqref{eq:whiteford}, trong đó độ đỏ và độ bão hòa làm tăng Arousal, còn độ sáng điều chỉnh mức nền. Hình~\ref{fig:color-flow} tóm tắt luồng ánh xạ này.
\begin{equation}\label{eq:whiteford}
    A = 0.373 \cdot \text{redness} + 0.356 \cdot \text{saturation} + 0.271 \cdot \text{lightness} - 0.10
\end{equation}

\begin{figure}[h!]
\centering
\begin{tikzpicture}[font=\small,
   blk/.style={draw, rounded corners, align=center, minimum height=0.8cm},
   arr/.style={-{Stealth[length=2mm]}, thick}]
  \node[blk, fill=black!4, text width=1.8cm] (hex) at (0,0) {Mã màu \texttt{\#RRGGBB}};
  \node[blk, fill=blue!6, text width=1.8cm] (oklab) at (2.8,0) {Oklab $(L,a,b)$};
  \node[blk, fill=blue!8, text width=2.6cm] (val) at (6.2,0.9) {Valence: ridge $\rightarrow$ ICEAS};
  \node[blk, fill=blue!8, text width=2.6cm] (aro) at (6.2,-0.9) {Arousal: Whiteford (Pt.~\ref{eq:whiteford})};
  \node[blk, fill=green!7, text width=1.8cm] (va) at (9.6,0) {Điểm V-A};
  \draw[arr] (hex)--(oklab);
  \draw[arr] (oklab) |- (val);
  \draw[arr] (oklab) |- (aro);
  \draw[arr] (val) -| (va);
  \draw[arr] (aro) -| (va);
\end{tikzpicture}
\caption{Luồng ánh xạ màu sắc sang tọa độ Valence--Arousal}
\label{fig:color-flow}
\end{figure}

Tính đúng đắn của công thức Arousal được kiểm chứng độc lập bằng nhịp độ thật của các bài được truy hồi: tương quan giữa độ sáng của màu với nhịp độ đạt 0.41 và giữa độ bão hòa với nhịp độ đạt 0.66, đúng theo hướng màu sáng và bão hòa hơn ứng với nhạc nhanh hơn.

\subsection{Truy hồi, chấm điểm và sắp xếp}
\label{subsection:4.2.3}
Với chức năng gợi ý theo màu, hệ thống chấm điểm độ gần cảm xúc bằng một hàm Gaussian RBF dị hướng (Phương trình~\eqref{eq:rbf}) với độ rộng trên trục Valence là $\sigma_V = 0.20$ và trên trục Arousal là $\sigma_A = 0.14$, do phân bố Arousal của kho nhạc nén hơn. Việc so khớp được thực hiện trong không gian phân vị (CDF) của kho nhạc (Phương trình~\eqref{eq:cdf}) để các màu khác nhau thực sự kéo về các vùng cảm xúc khác nhau.
\begin{equation}\label{eq:rbf}
    s(\mathbf{q}, \mathbf{p}) = \exp\!\left(-\frac{(q_V - p_V)^2}{2\sigma_V^2} - \frac{(q_A - p_A)^2}{2\sigma_A^2}\right)
\end{equation}
\begin{equation}\label{eq:cdf}
    \mathbf{q} = \big(F_V(c_V),\, F_A(c_A)\big), \qquad
    F_X(x) = \frac{1}{N}\big|\{\, i : X_i \le x \,\}\big|
\end{equation}
trong đó $\mathbf{q}$ là đích phân vị của màu (qua hàm phân phối tích lũy thực nghiệm $F$ của kho nhạc) và $\mathbf{p}$ là phân vị V-A của một bài. Để các bài trong cùng một màu nghe đồng nhất, hệ thống lấy dư một tập ứng viên rồi chọn cụm có độ nhất quán âm thanh cao nhất, đo bằng cosine trên biểu diễn MuQ đã được trừ trung bình để khử dị hướng \cite{ethayarajh2019anisotropy, mu2018allbutthetop}. Trọng số giữa độ khớp cảm xúc và độ nhất quán âm thanh được đặt ở mức 0.55. Việc đa dạng hóa theo nghệ sĩ \textit{không} dùng một bước MMR riêng mà được thực hiện ngay trong quá trình chọn cụm: mỗi khi một bài được chọn, các bài còn lại cùng nghệ sĩ bị phạt giảm điểm. Cuối cùng, danh sách được sắp xếp thành chuỗi chuyển cảm xúc mượt theo nguyên lý đồng cảm \cite{starcke2021iso}. Thuật toán~\ref{alg:color} trình bày toàn bộ quy trình.

\begin{algorithm}[h]
\caption{Gợi ý playlist theo màu}
\label{alg:color}
\KwIn{Một hoặc hai mã màu, số kết quả $k$}
\KwOut{Danh sách $k$ bài khớp cảm xúc và nhất quán âm thanh}
Ánh xạ mỗi màu sang điểm V-A (Hình~\ref{fig:color-flow})\;
\eIf{hai màu}{
  Sinh chuỗi điểm đích theo nguyên lý đồng cảm giữa hai màu (iso-principle)\;
}{
  Quy màu về đích phân vị $\mathbf{q}$ trong không gian CDF (Pt.~\eqref{eq:cdf})\;
}
Chấm điểm mọi bài bằng RBF dị hướng (Pt.~\eqref{eq:rbf}); lấy dư tập ứng viên $C$\;
Chọn tham lam cụm con của $C$ có độ nhất quán âm thanh cao nhất (cosine MuQ đã trừ trung bình, trọng số 0.55), đồng thời phạt giảm điểm các bài trùng nghệ sĩ ngay trong khi chọn để đa dạng hóa\;
\Return danh sách đã sắp theo trật tự cảm xúc mượt\;
\end{algorithm}

Với chức năng gợi ý bài tương tự, điểm tương đồng giữa hai bài là tổ hợp tuyến tính của ba thành phần như trong Phương trình~\eqref{eq:fusion}: tương đồng âm thanh trên biểu diễn MuQ chiếm trọng số 0.76, độ gần cảm xúc chiếm 0.16 và tương đồng nội dung lời chiếm 0.08. Các bản cover và trùng được loại bằng một danh sách dựng sẵn; sau đó kết quả được đa dạng hóa bằng MMR \cite{carbonell1998mmr} trên biểu diễn lời (e5-large) (hệ số $\lambda = 0.7$, cân bằng độ liên quan với bài nguồn và độ khác biệt giữa các kết quả). Vì các bài cùng một nghệ sĩ thường nằm sát nhau trong không gian biểu diễn, bước MMR này cũng \textit{gián tiếp} giảm hiện tượng một nghệ sĩ chiếm nhiều vị trí. Ngoài ra hệ thống hỗ trợ một tùy chọn vận hành giới hạn cứng số bài tối đa mỗi nghệ sĩ, nhưng tùy chọn này \textit{mặc định tắt} (Thuật toán~\ref{alg:similar}).
\begin{equation}\label{eq:fusion}
    \text{score} = 0.76 \cdot \cos(\text{MuQ}) + 0.16 \cdot \text{sim}(\text{V-A}) + 0.08 \cdot \cos(\text{lyrics})
\end{equation}

\begin{algorithm}[h]
\caption{Gợi ý bài hát tương tự}
\label{alg:similar}
\KwIn{Chỉ số bài nguồn $s$, số kết quả $k$, tập loại trừ $E$ (chế độ đài)}
\KwOut{Danh sách $k$ bài tương tự}
\ForEach{bài $i \neq s$, $i \notin E$}{
  Tính $\text{score}(s, i)$ theo Pt.~\eqref{eq:fusion}\;
}
Loại các bản cover/trùng của $s$ theo chỉ mục cover dựng sẵn\;
Đa dạng hóa bằng MMR trên biểu diễn lời (e5-large), $\lambda = 0.7$ (Pt.~\eqref{eq:mmr}) --- gián tiếp giảm trùng nghệ sĩ\;
\textit{(Tùy chọn, mặc định tắt)} giới hạn cứng số bài tối đa mỗi nghệ sĩ\;
\Return $k$ bài điểm cao nhất\;
\end{algorithm}

Hai sơ đồ tuần tự trong Hình~\ref{fig:seq-uc01} và Hình~\ref{fig:seq-uc02} mô tả tương tác giữa các thành phần khi thực thi hai chức năng cốt lõi.

\begin{figure}[h!]
\centering
\begin{tikzpicture}[font=\footnotesize,
   hd/.style={draw, rounded corners, fill=black!4, align=center, minimum width=2.2cm, minimum height=0.6cm},
   msg/.style={-{Stealth[length=1.8mm]}, thick}]
  \def\ya{0} \def\yb{-4.9}
  \foreach \x/\n in {0/Người nghe, 2.9/SPA, 5.8/API, 8.4/Engine} {
    \node[hd] (h\n) at (\x,\ya) {\n};
    \draw[dashed, gray] (\x,\ya-0.35) -- (\x,\yb);
  }
  \draw[msg] (0,-0.9) -- node[above]{chọn bài hát} (2.9,-0.9);
  \draw[msg] (2.9,-1.5) -- node[above]{\texttt{GET /song/\{id\}/similar}} (5.8,-1.5);
  \draw[msg] (5.8,-2.1) -- node[above]{\texttt{recommend\_by\_song()}} (8.4,-2.1);
  \draw[msg] (8.4,-2.7) to[out=-35,in=35,looseness=3.5] node[right, align=left, text width=2.4cm, font=\scriptsize]{tổ hợp + loại cover + MMR} (8.4,-3.3);
  \draw[msg] (8.4,-3.7) -- node[above]{danh sách top-$k$} (5.8,-3.7);
  \draw[msg] (5.8,-4.1) -- node[above]{JSON} (2.9,-4.1);
  \draw[msg] (2.9,-4.5) -- node[above]{hiển thị} (0,-4.5);
\end{tikzpicture}
\caption{Sơ đồ tuần tự UC01 --- gợi ý bài tương tự}
\label{fig:seq-uc01}
\end{figure}

\begin{figure}[h!]
\centering
\begin{tikzpicture}[font=\footnotesize,
   hd/.style={draw, rounded corners, fill=black!4, align=center, minimum width=2.2cm, minimum height=0.6cm},
   msg/.style={-{Stealth[length=1.8mm]}, thick}]
  \def\ya{0} \def\yb{-4.9}
  \foreach \x/\n in {0/Người nghe, 2.9/SPA, 5.8/API, 8.4/Engine} {
    \node[hd] (g\n) at (\x,\ya) {\n};
    \draw[dashed, gray] (\x,\ya-0.35) -- (\x,\yb);
  }
  \draw[msg] (0,-0.9) -- node[above]{chọn 1--2 màu} (2.9,-0.9);
  \draw[msg] (2.9,-1.5) -- node[above]{\texttt{POST /recommend/color}} (5.8,-1.5);
  \draw[msg] (5.8,-2.1) -- node[above]{\texttt{recommend\_by\_colors()}} (8.4,-2.1);
  \draw[msg] (8.4,-2.7) to[out=-35,in=35,looseness=3.5] node[right, align=left, text width=2.4cm, font=\scriptsize]{map $\rightarrow$ CDF $\rightarrow$ RBF $\rightarrow$ cụm nhất quán (phạt nghệ sĩ)} (8.4,-3.3);
  \draw[msg] (8.4,-3.7) -- node[above]{playlist + \texttt{why}/\texttt{bridge}} (5.8,-3.7);
  \draw[msg] (5.8,-4.1) -- node[above]{JSON} (2.9,-4.1);
  \draw[msg] (2.9,-4.5) -- node[above]{hiển thị} (0,-4.5);
\end{tikzpicture}
\caption{Sơ đồ tuần tự UC02 --- gợi ý theo màu}
\label{fig:seq-uc02}
\end{figure}

\subsection{Ví dụ minh họa}
\label{subsection:4.2.4}
Để làm rõ cách hai chức năng vận hành, mục này trình bày hai ví dụ với số liệu thật từ hệ thống. Xét truy vấn gợi ý theo \textbf{màu lam} (\texttt{\#0067A5}). Hệ thống ánh xạ màu này sang tọa độ cảm xúc $(V, A) = (0.565,\ 0.358)$: Valence trên mức trung bình và Arousal dưới mức trung bình, tức rơi vào góc phần tư ``thư thái / bình yên'' (Q4). Tọa độ này được quy về đích phân vị trong phân bố kho nhạc, rồi hàm RBF dị hướng chấm điểm mọi bài theo độ gần với đích; hệ thống lấy dư một tập ứng viên, chọn cụm có độ nhất quán âm thanh cao nhất trên biểu diễn MuQ đã trừ trung bình, và đa dạng hóa theo nghệ sĩ. Kết quả là một playlist gồm các bài nhẹ nhàng, năng lượng thấp, nghe đồng nhất với nhau. Ngược lại, nếu người dùng chọn \textbf{màu vàng} (\texttt{\#F3C300}) với $(V, A) = (0.767,\ 0.696)$ --- Valence và Arousal đều cao, thuộc góc ``vui vẻ / phấn khích'' (Q1) --- hệ thống trả về các bài sôi động, vui tươi. Sự tương phản rõ rệt giữa hai playlist này chính là biểu hiện trực quan của độ tách biệt giữa các màu mà mục~\ref{subsection:eval-color} đo định lượng.

Với chức năng gợi ý bài tương tự, giả sử bài nguồn là một bản ballad có tọa độ cảm xúc gần $(V, A) = (0.54,\ 0.46)$. Hệ thống tính điểm tổ hợp (Phương trình~\eqref{eq:fusion}) giữa bài nguồn và mọi bài khác: trọng số lớn nhất dành cho tương đồng âm thanh trên MuQ (nên kết quả ``nghe giống'' về hòa âm, nhịp điệu), bổ sung bởi độ gần cảm xúc và tương đồng lời. Trước khi trả về, các bản cover/trùng của bài nguồn bị loại theo chỉ mục dựng sẵn (tránh trả về chính bài nguồn dưới một bản phối khác), và bước đa dạng hóa MMR trên biểu diễn lời (e5-large) giúp danh sách không bị một nghệ sĩ chiếm trọn (vì các bài cùng nghệ sĩ nằm gần nhau trong không gian biểu diễn nên bị MMR trải ra). Nếu người nghe tiếp tục yêu cầu (chế độ ``đài''), tập loại trừ $E$ được nối thêm các bài đã phát để danh sách không lặp lại.

\section{Công cụ và mô hình sử dụng}
\label{section:4.3}
Bảng~\ref{table:tools} liệt kê các công cụ và mô hình chính được sử dụng. Toàn bộ mô hình học máy được dùng ở dạng đóng băng; phần huấn luyện duy nhất là các đầu dò tuyến tính (hồi quy ridge) trên các tập dữ liệu công khai. Phần phục vụ được hiện thực bằng FastAPI; giao diện là một ứng dụng trang đơn React. Tầng dữ liệu dùng PostgreSQL~17 kèm pgvector và pg\_trgm (truy cập qua SQLAlchemy 2.0, lược đồ quản lý bằng Alembic) để lưu siêu dữ liệu, hạt giống và dữ liệu phụ cho crossfade, cùng tìm kiếm mờ tiếng Việt. Tuy nhiên cần nói rõ: \textit{đường phục vụ truy vấn là file-based} --- các đặc trưng tính sẵn được nạp từ tệp \texttt{.npy}/\texttt{.json} vào bộ nhớ, nên cơ sở dữ liệu \textit{không} nằm trên đường truy vấn nóng (chỉ nạp một lần lúc khởi động, có cơ chế lùi an toàn nếu thiếu). Toàn hệ thống được đóng gói bằng Docker; pha thu thập dữ liệu ngoại tuyến dùng ytmusicapi, yt-dlp và LRCLIB.

\begin{table}[h!]
\centering
\caption{Danh sách công cụ và mô hình chính}
\label{table:tools}
\begin{tabular}{p{4.2cm} p{4.0cm} p{3.2cm}}
\toprule
\textbf{Mục đích} & \textbf{Công cụ / mô hình} & \textbf{Nguồn} \\
\midrule
Ngôn ngữ và xử lý dữ liệu & Python, NumPy, scikit-learn & --- \\
Biểu diễn âm nhạc & MuQ & \cite{zhu2025muq} \\
Cảm xúc lời (ngữ cảnh VN) & ViSoBERT & \cite{nguyen2023visobert} \\
Cảm xúc lời (liên ngữ) & XLM-R & \cite{conneau2020xlmr} \\
Tương đồng lời & multilingual-e5-large & \cite{wang2024e5} \\
Từ điển Valence & NRC-VAD, VnEmoLex & \cite{mohammad2018nrcvad, doan2022vnemolex} \\
Phục vụ (API) & FastAPI (Python) & --- \\
Giao diện & Ứng dụng trang đơn React & --- \\
Tầng dữ liệu \& tìm kiếm & PostgreSQL 17 + pgvector + pg\_trgm & --- \\
ORM \& di trú lược đồ & SQLAlchemy 2.0, Alembic & --- \\
Đóng gói \& triển khai & Docker (compose) & --- \\
Thu thập dữ liệu & ytmusicapi, yt-dlp, LRCLIB & --- \\
\bottomrule
\end{tabular}
\end{table}

\section{Dữ liệu và kho nhạc}
\label{section:data}
Kho nhạc phục vụ gồm \textbf{5.138 bài hát nhạc Việt}, mỗi bài có tệp âm thanh, siêu dữ liệu (tên bài, nghệ sĩ, album) và lời bài hát. Đây là số bài thực tế mà động cơ gợi ý nạp lên và là cỡ tập đánh giá trong các thực nghiệm ở mục~\ref{section:eval}.

\textbf{Nguồn dữ liệu.} Dữ liệu được thu thập và xử lý tự động qua một \textit{quy trình bảy pha có điểm kiểm tra} (checkpoint) --- thu thập, lọc, tải, lấy lời, trích đặc trưng, xử lý nhãn, nạp kho --- cho phép chạy lại từng pha độc lập và tiếp tục khi gián đoạn mà không làm lại từ đầu, hoàn toàn không sử dụng dữ liệu người dùng (Hình~\ref{fig:collect}): (i) khám phá nghệ sĩ Việt từ siêu dữ liệu Spotify, khởi đầu từ 353 nghệ sĩ hạt giống; (ii) lấy danh sách bài hát, siêu dữ liệu và lời từ YouTube Music; (iii) tải tệp âm thanh bằng \texttt{yt-dlp} theo bốn mức dự phòng (ưu tiên bản nhạc khớp thời lượng) kèm cắt khoảng lặng đầu/cuối; (iv) bổ sung lời bài hát từ dịch vụ LRCLIB. Như vậy Spotify chỉ phục vụ khám phá/định danh nghệ sĩ, còn âm thanh và lời lấy từ YouTube Music và LRCLIB.

\begin{figure}[h!]
\centering
\begin{tikzpicture}[font=\footnotesize,
   stg/.style={draw, rounded corners, align=center, fill=black!3, text width=2.3cm, minimum height=0.85cm},
   arr/.style={-{Stealth[length=2mm]}, thick}]
  \node[stg] (seed) at (0,0) {353 nghệ sĩ hạt giống (Spotify)};
  \node[stg] (yt) at (3.0,0) {YT Music: bài + lời + siêu dữ liệu};
  \node[stg] (dl) at (6.0,0) {\texttt{yt-dlp} tải âm thanh (4 mức dự phòng)};
  \node[stg] (lrc) at (9.0,0) {LRCLIB bổ sung lời};
  \node[stg, fill=green!7] (out) at (11.7,0) {5.138 bài + chỉ mục cover};
  \draw[arr] (seed)--(yt); \draw[arr] (yt)--(dl); \draw[arr] (dl)--(lrc); \draw[arr] (lrc)--(out);
\end{tikzpicture}
\caption{Quy trình thu thập dữ liệu tự động, không dùng dữ liệu người dùng}
\label{fig:collect}
\end{figure}

\textbf{Tiền xử lý âm thanh.} Trước khi trích embedding, mỗi tệp âm thanh được cắt khoảng lặng ở đầu và cuối để tránh nhiễu cho biểu diễn; việc tải ưu tiên bản nhạc có thời lượng khớp với siêu dữ liệu nhằm loại các bản remix hoặc bản kéo dài bất thường. Nhịp độ được trích bằng phương pháp dò phách (downbeat) của thư viện librosa thay vì dùng giá trị nhịp độ thô, vì giá trị thô dễ bị nhân đôi hoặc chia đôi (nhầm 60 BPM thành 120 BPM); nhịp độ ``sạch'' này quan trọng vì nó là một trong ba thành phần của trục Arousal (mục~\ref{subsection:4.2.1}). Các đặc trưng âm thanh từ mô hình Essentia ở tần số lấy mẫu cao bị suy biến nên không được dùng; chỉ giữ lại các đặc trưng tín hiệu số đáng tin như nhịp độ và độ to.

\textbf{Lời bài hát và xử lý bài thiếu lời.} Trên kho hiện tại, \textbf{toàn bộ 5.138/5.138 bài (100\%) đều có lời} (cờ \texttt{has\_lyrics}). Hệ thống vẫn có cơ chế dự phòng cho trường hợp thiếu lời: khi một bài không có lời, tín hiệu Valence từ lời không khả dụng nên tọa độ cảm xúc được suy từ các đặc trưng âm thanh (valence/\allowbreak arousal/\allowbreak energy) thay thế. Vì kho hiện tại phủ lời 100\%, cơ chế này không phải kích hoạt trong thực tế nhưng vẫn bảo đảm hệ thống không gãy khi mở rộng kho.

\textbf{Xử lý bài trùng và cover.} Các phiên bản trùng/cover được phát hiện ở pha ngoại tuyến và lưu thành một chỉ mục dựng sẵn (\texttt{cover\_index.json}). Việc phát hiện lấy \textit{lời làm tín hiệu chính} (độ tương đồng cosine của lời rất cao tương ứng cùng nội dung) kết hợp khớp tiêu đề đã chuẩn hóa (bắt được cả phiên bản song ngữ cùng bài). Kết quả hiện tại: \textbf{1.124 bài} có quan hệ cover, gom thành \textbf{979 cặp}. Khi gợi ý bài tương tự, các cover của bài nguồn bị loại để không lãng phí vị trí và tránh trả về ``cùng một bài''.

\textbf{Phân bố cảm xúc của kho nhạc.} Hình~\ref{fig:scatter} biểu diễn phân bố Valence--Arousal của toàn bộ 5.138 bài, còn Hình~\ref{fig:hist} cho thấy phân bố của từng trục. Có thể thấy trục Arousal nén hơn trục Valence (độ lệch chuẩn $\approx 0.13$ so với $\approx 0.18$); đây chính là lý do hàm RBF dùng độ rộng nhỏ hơn trên trục Arousal ($\sigma_A = 0.14 < \sigma_V = 0.20$). Vùng năng lượng cao (Arousal lớn) thưa hơn, một hạn chế được nêu lại ở Chương~\ref{chapter:conclusion}.

\begin{figure}[h!]
\centering
\begin{subfigure}{0.48\textwidth}
\centering
\includegraphics[width=\textwidth]{Hinhve/fig_va_scatter.png}
\caption{Phân bố V-A (5.138 bài)}
\label{fig:scatter}
\end{subfigure}\hfill
\begin{subfigure}{0.5\textwidth}
\centering
\includegraphics[width=\textwidth]{Hinhve/fig_va_hist.png}
\caption{Histogram Valence và Arousal}
\label{fig:hist}
\end{subfigure}
\caption{Phân bố cảm xúc của kho nhạc: trục Arousal nén hơn trục Valence}
\label{fig:va-dist}
\end{figure}

\textbf{Lưu trữ đặc trưng.} Các đặc trưng nặng được tính sẵn và lưu ra tệp để pha phục vụ nạp vào bộ nhớ (Bảng~\ref{table:features}). Đặc trưng MuQ được lưu ở \texttt{muq\_embeddings.npy} dưới dạng ma trận $5138 \times 1024$ (số thực 32-bit, đã chuẩn hóa L2) kèm tệp siêu dữ liệu căn theo thứ tự kho nhạc. Lược đồ đầy đủ được trình bày trong Phụ lục~\ref{appendix:schema}.

\begin{table}[h!]
\centering
\caption{Các tệp đặc trưng tính sẵn được pha phục vụ nạp lên}
\label{table:features}
\begin{tabular}{p{4.2cm} p{5.4cm} p{2.8cm}}
\toprule
\textbf{Đặc trưng} & \textbf{Tệp lưu trữ} & \textbf{Kích thước} \\
\midrule
Embedding âm thanh MuQ & \texttt{muq\_embeddings.npy} & $5138 \times 1024$ \\
Embedding nội dung lời & \texttt{lyrics\_e5large.npy} & $5138 \times 1024$ \\
Tọa độ cảm xúc (V-A) & \texttt{emotion\_\allowbreak labels\_\allowbreak v6i.json} & 5138 mục \\
Nhịp độ sạch (BPM) & \texttt{clean\_bpm.json} & theo bài \\
Chỉ mục cover/trùng & \texttt{cover\_index.json} & 1124 bài \\
\bottomrule
\end{tabular}
\end{table}

\section{Giao diện lập trình ứng dụng}
\label{section:api}
Tầng phục vụ được hiện thực bằng FastAPI và cung cấp một tập endpoint cho ứng dụng trang đơn (Bảng~\ref{table:api}). Các phản hồi dùng một phong bì JSON nhất quán với trường \texttt{success} và phần dữ liệu (\texttt{results}/\texttt{songs}/\texttt{song}), kèm \texttt{message} khi cần.

\begin{table}[h!]
\centering
\caption{Các endpoint cốt lõi của giao diện lập trình ứng dụng (toàn hệ thống có khoảng 22 endpoint)}
\label{table:api}
\begin{tabular}{p{1.8cm} p{5.8cm} p{6.0cm}}
\toprule
\textbf{Phương thức} & \textbf{Đường dẫn} & \textbf{Chức năng} \\
\midrule
POST & \texttt{/api/recommend/color} & Gợi ý playlist theo 1--2 màu (UC02) \\
GET & \texttt{/api/song/\{id\}/similar} & Gợi ý bài tương tự (UC01) \\
GET & \texttt{/api/song/\{id\}} & Chi tiết + thông tin cảm xúc của bài (UC04) \\
GET & \texttt{/api/songs} & Duyệt/lọc/sắp xếp kho nhạc (UC03) \\
GET & \texttt{/api/search} & Tìm kiếm theo tên/nghệ sĩ/lời/cảm giác \\
GET & \texttt{/api/audio/stream/\{id\}} & Phát tệp âm thanh (hoặc chuyển hướng CDN) \\
\bottomrule
\end{tabular}
\end{table}

Mỗi bài trong kết quả gồm: chỉ số (\texttt{song\_index}), định danh (\texttt{track\_id}), tên bài, nghệ sĩ, mã màu tâm trạng (\texttt{color\_hex}), các giá trị Valence/Arousal/Energy và nhãn góc phần tư cảm xúc (\texttt{mood\_quadrant}). Với gợi ý theo màu, mỗi bài kèm một đối tượng giải thích \texttt{why} gồm: câu giải thích tiếng Việt (\texttt{reason}), tín hiệu mạnh nhất (\texttt{top\_signal}), điểm khớp cảm xúc (\texttt{mood\_match}) và tọa độ V-A của bài (\texttt{song\_va}) so với của màu (\texttt{color\_va}). Với gợi ý bài tương tự, mỗi bài kèm điểm \texttt{similarity\_score} cùng cờ và đường dẫn phát âm thanh. Quan hệ màu$\rightarrow$cảm xúc tổng quát được trả riêng ở khối siêu dữ liệu \texttt{query.bridge} (và \texttt{query.journey} khi chọn hai màu), \textit{không} nằm trong từng \texttt{why}. Một ví dụ đầy đủ về yêu cầu và phản hồi JSON được trình bày trong Phụ lục~\ref{appendix:api}.

\textbf{Xử lý lỗi ở biên.} Hệ thống xác thực đầu vào và xử lý các trường hợp biên một cách tường minh: mã màu không hợp lệ (không đúng dạng \texttt{\#RRGGBB}) bị từ chối bằng lỗi xác thực đầu vào; bài hát không tồn tại trả về mã 404; bài thiếu lời thì trường lời trả về rỗng (không phải lỗi); bài thiếu một đặc trưng âm thanh thì trường tương ứng trả về rỗng thay vì gây lỗi; định danh sai định dạng khi phát trả về mã 400 và tệp âm thanh không có trả về 404. Các endpoint gợi ý có giới hạn tần suất để tránh quá tải.

\section{Giao diện người dùng}
\label{section:ui}
Giao diện là một ứng dụng trang đơn React, cung cấp \textbf{hai giao diện (skin) dùng chung một logic gợi ý}: (i) giao diện ``đắm chìm'' 3D mô phỏng hệ mặt trời, trong đó mỗi màu là một hành tinh và hành trình hai màu là một chuyến bay; và (ii) giao diện ``classic'' 2D theo phong cách trình phát nhạc quen thuộc. Người dùng đổi qua lại giữa hai giao diện và lựa chọn được lưu cục bộ.

Các màn hình chính gồm: màn chọn màu (bảng 12 màu/hành tinh), màn kết quả playlist theo màu, màn duyệt và chọn bài hát, màn kết quả bài tương tự (chế độ ``đài''), bảng lời bài hát, thanh phát nhạc, ô tìm kiếm và lớp hướng dẫn sử dụng. Tâm trạng (Valence--Arousal) của mỗi bài được thể hiện trực quan qua màu nền ảnh bìa và một nhãn tâm trạng ngắn hiện khi đổi bài. Hai giao diện này không phải hai ứng dụng tách biệt mà là hai ``lớp da'' (skin) khác nhau trên cùng một logic gợi ý và cùng một tầng dữ liệu: mọi lời gọi API, mọi kết quả và cách xử lý đều dùng chung, chỉ khác ở cách trình bày. Giao diện 3D phù hợp với trải nghiệm khám phá giàu cảm xúc (mỗi màu là một hành tinh, hành trình hai màu là một chuyến bay), trong khi giao diện 2D quen thuộc với người dùng đã quen các trình phát nhạc thông thường. Việc chia sẻ logic giữa hai giao diện vừa giảm trùng lặp mã nguồn vừa bảo đảm hành vi gợi ý nhất quán bất kể người dùng chọn giao diện nào.

Ngoài ra, trình phát hỗ trợ chuyển bài mượt (crossfade) dựa trên các điểm cue và độ to (LUFS) đã tính sẵn: hệ thống bắt đầu làm mờ bài đang phát và đưa bài kế tiếp vào tại các điểm chuyển hợp lý, đồng thời cân bằng độ to giữa hai bài để tránh chênh lệch âm lượng đột ngột. Đây là tính năng phát nhạc bổ trợ, không nằm trong hai chức năng cốt lõi của đồ án.

Cần nêu rõ một điểm trung thực về mức hoàn thiện: dữ liệu cảm xúc đã được tính và \textit{cung cấp đầy đủ qua API} (tọa độ V-A, nhãn tâm trạng, đối tượng \texttt{why}, cầu nối màu$\rightarrow$cảm xúc), nhưng một \textit{màn chi tiết hiển thị tọa độ Valence--Arousal dạng số} và một \textit{bảng giải thích ``vì sao bài được gợi ý''} chưa được dựng thành màn riêng trong giao diện hiện tại; chúng được xem là hướng hoàn thiện trình bày ở Chương~\ref{chapter:conclusion}. Các màn hình chính được minh họa lần lượt như sau: với chức năng gợi ý theo màu là màn chọn màu (Hình~\ref{fig:ui-color}) và màn kết quả playlist theo màu với ảnh bìa nhuộm theo tâm trạng (Hình~\ref{fig:ui-color-results}); với chức năng gợi ý bài tương tự là màn duyệt/chọn bài hát (Hình~\ref{fig:ui-browse}) và màn kết quả bài tương tự ở chế độ ``đài'' nối tiếp (Hình~\ref{fig:ui-similar}); ngoài ra là màn thông tin cảm xúc của bài hát (Hình~\ref{fig:ui-detail}) và phần giải thích ngắn ``vì sao bài được gợi ý'' (Hình~\ref{fig:ui-why}). Do hệ thống chạy tương tác (đặc biệt là giao diện 3D), ảnh chụp màn hình thực tế được chèn ở các vị trí tương ứng.

\begin{figure}[h!]
\centering
\uiscreenshot{Màn chọn màu (bảng 12 màu / hành tinh)}
\caption{Màn chọn màu để gợi ý playlist (UC02)}
\label{fig:ui-color}
\end{figure}

\begin{figure}[h!]
\centering
\uiscreenshot{Kết quả playlist theo màu --- danh sách bài kèm ảnh bìa nhuộm theo tâm trạng}
\caption{Màn kết quả playlist theo màu (UC02, UC03)}
\label{fig:ui-color-results}
\end{figure}

\begin{figure}[h!]
\centering
\uiscreenshot{Màn duyệt/chọn bài hát (thư viện, bài nổi bật, mới phát hành)}
\caption{Màn chọn bài hát để tìm bài tương tự (UC01)}
\label{fig:ui-browse}
\end{figure}

\begin{figure}[h!]
\centering
\uiscreenshot{Kết quả bài tương tự --- hàng đợi ``đài'' nối tiếp}
\caption{Màn kết quả bài tương tự (UC01, UC03)}
\label{fig:ui-similar}
\end{figure}

\begin{figure}[h!]
\centering
\uiscreenshot{Thông tin cảm xúc của bài hát --- màu nền theo V-A, nhãn tâm trạng, lời bài hát (màn V-A dạng số là hướng hoàn thiện)}
\caption{Màn thông tin cảm xúc của bài hát (UC04)}
\label{fig:ui-detail}
\end{figure}

\begin{figure}[h!]
\centering
\uiscreenshot{Giải thích ngắn vì sao bài được gợi ý (đối tượng \texttt{why} từ API --- hiển thị trong giao diện là hướng hoàn thiện)}
\caption{Giải thích lý do gợi ý dựa trên cảm xúc}
\label{fig:ui-why}
\end{figure}

\section{Đánh giá thực nghiệm}
\label{section:eval}
Do không có dữ liệu người dùng, hệ thống được đánh giá bằng nhiều tầng độ đo ngoại tuyến không vòng tròn, bổ sung bằng các kiểm chứng độc lập (nhịp độ thật, chuyển miền sang tập khác). Nguyên tắc xuyên suốt là độ đo cuối quyết định việc chấp nhận một thành phần, và mọi thay đổi phải vượt một cổng kiểm định không thụt lùi.

\subsection{Phương pháp và dữ liệu đánh giá}
\label{subsection:eval-method}
Tập đánh giá là toàn bộ \textbf{5.138 bài hát} của kho nhạc. Chức năng gợi ý theo màu được đánh giá trên \textbf{12 màu cơ bản} (theo Berlin \& Kay \cite{berlin1969basic}, cũng là 12 màu mà ICEAS \cite{jonauskaite2020universal} khảo sát) với mỗi truy vấn lấy $top\text{-}k=10$.

Vì không có tập nhạc Việt được con người gán nhãn màu--cảm xúc, \textbf{ground truth cho gợi ý theo màu không phải nhãn người} mà được tạo theo phương pháp hiệu chỉnh phân bố kho nhạc \cite{steck2018calibrated}: mỗi màu được ánh xạ thành một tọa độ Valence--Arousal, rồi quy về \textit{đích phân vị} tương ứng trong phân bố tích lũy (CDF) của kho nhạc. Độ đo chính là \textit{sai số nhắm trúng đích} (targeting error, TE) theo Phương trình~\eqref{eq:te}: khoảng cách Euclid trong không gian phân vị giữa đích phân vị của màu truy vấn và phân vị V-A của các bài được gợi ý, lấy trung bình trên các bài rồi trung bình trên 12 màu.
\begin{equation}\label{eq:te}
    \text{TE} = \frac{1}{|C|}\sum_{c \in C} \frac{1}{k}\sum_{s \in \text{top-}k(c)} \big\lVert \mathbf{q}_c - \mathbf{p}_s \big\rVert_2
\end{equation}
trong đó $C$ là tập 12 màu, $\mathbf{q}_c$ là đích phân vị của màu $c$ và $\mathbf{p}_s$ là phân vị V-A của bài $s$. Đây là một độ đo ngoại tuyến và \textit{tự tham chiếu} (đo trên chính nhãn V-A dùng để khớp); hạn chế này được nêu rõ ở cuối mục. Để tránh kết luận chỉ dựa vào một độ đo tự tham chiếu, TE được bổ sung bằng các kiểm chứng độc lập với nhãn (nhịp độ thật, chuẩn ICEAS, chuyển miền PMEmo) trình bày bên dưới. Các số liệu trong mục này được sinh bởi bộ công cụ đánh giá ngoại tuyến (\texttt{tools/color\_eval\_rigor.py}, \texttt{tools/color\_ablation.py}) và có thể chạy lại để tái lập.

\textbf{Phương pháp so sánh và lựa chọn mô hình.} Với mỗi vị trí có thể thay mô hình (trục biểu diễn âm thanh, embedding lời, các đầu dò cảm xúc, mô hình màu--cảm xúc), việc lựa chọn không dựa vào thứ hạng trên các bộ chuẩn quốc tế --- vì ``tốt trên giấy'' không bảo đảm tốt cho nhiệm vụ và dữ liệu cụ thể của hệ thống --- mà tuân theo năm nguyên tắc. Thứ nhất, (i) phán quyết bằng \textit{độ đo cuối} gắn với chính nhiệm vụ: NDCG trên ground truth tương đồng âm nhạc cho gợi ý bài tương tự, TE cùng các kiểm chứng độc lập cho gợi ý theo màu, và $R^2$ kiểm chứng chéo trên DEAM cho các đầu dò cảm xúc. Thứ hai, (ii) \textit{so sánh có kiểm soát}: chỉ thay đúng thành phần cần so, giữ nguyên mọi phần còn lại. Thứ ba, (iii) \textit{tinh chỉnh lại siêu tham số riêng} cho từng ứng viên trước khi so, vì kế thừa thiết lập của mô hình cũ có thể khiến một mô hình thực sự tốt hơn lại bị đánh giá là kém. Thứ tư, (iv) bảo đảm tính vững và chống hiện tượng ``tối ưu hóa chính độ đo'' bằng kiểm chứng chéo $k$-fold (tinh chỉnh trên fold huấn luyện, chấm trên fold tách riêng), khoảng tin cậy bootstrap và kiểm định ý nghĩa (Wilcoxon kèm hiệu chỉnh FDR), kèm \textit{tam giác hóa} bằng một tín hiệu độc lập (chuyển miền sang tập khác hoặc một bộ chấm độc lập). Cuối cùng, (v) chỉ chấp nhận thay thế khi ứng viên \textit{thắng hoặc hòa} trên độ đo cuối và vượt \textit{cổng không thụt lùi}, ưu tiên mô hình gọn/nhẹ hơn khi hòa. Kết quả áp dụng quy trình này cho toàn bộ các thành phần được trình bày ở mục~\ref{subsection:eval-bakeoff}.

\subsection{Đánh giá chức năng gợi ý theo màu}
\label{subsection:eval-color}
Bằng chứng \textit{chính} cho chức năng này là các kiểm chứng độc lập với nhãn (khớp chuẩn người ICEAS và nhịp độ thật, trình bày bên dưới); sai số nhắm đích (TE) --- vì tự tham chiếu --- chỉ đóng vai trò một \textit{điều kiện cần}. Cấu hình hiện tại đạt \textbf{TE $= 0.0242$} với khoảng tin cậy bootstrap 95\% là $[0.0219,\,0.0263]$. Bảng~\ref{table:baseline} và Hình~\ref{fig:te} so sánh với năm baseline. Hệ thống vượt cả bốn baseline thực (ngẫu nhiên, theo độ phổ biến, chỉ Valence, chỉ Arousal): kiểm định Wilcoxon kèm hiệu chỉnh tỉ lệ phát hiện sai Benjamini--Hochberg \cite{benjamini1995fdr} cho \textbf{4/4 bác bỏ giả thuyết không với $p$ hiệu chỉnh $= 0.00031$}. Hệ thống không vượt baseline ``lân cận V-A tham lam'' --- đây là một \textit{trần lý thuyết} (oracle) loại bỏ đa dạng nên không phải đối thủ công bằng, và việc bám sát nó cho thấy phần đa dạng hóa chỉ đánh đổi rất ít TE.

\begin{table}[h!]
\centering
\caption{Sai số nhắm đích (TE) của hệ thống so với các baseline}
\label{table:baseline}
\begin{tabular}{p{6.6cm} p{3.0cm} p{3.2cm}}
\toprule
\textbf{Phương pháp} & \textbf{TE} & \textbf{Ghi chú} \\
\midrule
Ngẫu nhiên & 0.561 & baseline \\
Theo độ phổ biến & 0.513 & baseline \\
Chỉ Valence & 0.438 & loại bỏ một trục \\
Chỉ Arousal & 0.292 & loại bỏ một trục \\
Lân cận V-A tham lam (oracle) & 0.021 & trần lý thuyết, không đa dạng \\
\textbf{Brightify (V-A RBF)} & \textbf{0.0242} & CI95 $[0.0219, 0.0263]$ \\
\bottomrule
\end{tabular}
\end{table}

\begin{figure}[h!]
\centering
\includegraphics[width=0.78\textwidth]{Hinhve/fig_te_baselines.png}
\caption{So sánh TE của Brightify với các baseline (thấp hơn là tốt hơn); Brightify bám sát trần lý thuyết oracle}
\label{fig:te}
\end{figure}

Các kiểm chứng độc lập đều theo đúng hướng kỳ vọng. Valence của màu khớp chuẩn ICEAS với hệ số tương quan \textbf{$0.969$} (khoảng tin cậy Fisher-z $[0.889, 0.991]$); do chỉ có $n=12$ màu nên khoảng tin cậy rộng, và kiểm chứng giữ-một-ra (LOO-CV) cho $r \approx 0.87$. Công thức Arousal của Whiteford được kiểm bằng nhịp độ thật: $\rho(\text{độ sáng}, \text{BPM}) = 0.41$ và $\rho(\text{độ bão hòa}, \text{BPM}) = 0.66$. Hình~\ref{fig:colors} cho thấy 12 màu được ánh xạ vào các vùng cảm xúc khác nhau trên mặt phẳng V-A. Bảng~\ref{table:ablation} và Hình~\ref{fig:ablation} cho thấy đóng góp của hai cải tiến chính: so khớp trong không gian phân vị nâng \textbf{độ tách biệt giữa các màu từ $0.13$ lên $0.31$}, còn việc trừ trung bình khi đo nhất quán âm thanh nâng \textbf{độ nhất quán từ khoảng $0.02$ lên $0.23$}. Với hành trình hai màu, kiểm định Kolmogorov--Smirnov so với phân bố đều cho $\text{KS} = 0.21$ (đạt) và độ chênh nhịp độ giữa các bước giảm từ $31$ xuống $2.7$ BPM.

\begin{figure}[h!]
\centering
\begin{subfigure}{0.5\textwidth}
\centering
\includegraphics[width=\textwidth]{Hinhve/fig_colors_va.png}
\caption{12 màu trên mặt phẳng V-A}
\label{fig:colors}
\end{subfigure}\hfill
\begin{subfigure}{0.48\textwidth}
\centering
\includegraphics[width=\textwidth]{Hinhve/fig_ablation.png}
\caption{Đóng góp tích lũy (ablation)}
\label{fig:ablation}
\end{subfigure}
\caption{Trái: 12 màu kéo về các vùng cảm xúc khác nhau; Phải: mỗi cải tiến nâng độ tách biệt và độ nhất quán}
\label{fig:color-eval}
\end{figure}

\begin{table}[h!]
\centering
\caption{Ablation: mỗi cải tiến đóng góp vào tách biệt và nhất quán (12 màu, $top\text{-}k=10$)}
\label{table:ablation}
\begin{tabular}{p{5.4cm} c c c}
\toprule
\textbf{Cấu hình (cộng dồn)} & \textbf{Tách biệt} & \textbf{Nhất quán} & \textbf{TE} \\
\midrule
So khớp phân vị (cơ sở) & 0.129 & 0.015 & 0.0217 \\
+ Đích theo CDF & 0.225 & 0.026 & 0.0229 \\
+ Nhất quán âm thanh (trừ trung bình) & 0.224 & 0.113 & 0.0213 \\
+ Arousal Whiteford & 0.314 & 0.228 & 0.0242 \\
\bottomrule
\end{tabular}
\end{table}

Các kết quả này có một số hạn chế cần nêu thẳng: (i) TE là độ đo \textit{tự tham chiếu} (offline), chưa có nghiên cứu nghe thực tế của người dùng; (ii) ánh xạ màu--cảm xúc là \textit{ngoại suy}: liên kết màu--cảm xúc tuy có khuôn mẫu phổ quát nhưng vẫn mang tính đặc thù văn hóa \cite{jonauskaite2019ml}, trong khi Việt Nam không nằm trong 30 quốc gia của khảo sát ICEAS; trục Valence--Arousal được ghi nhận ổn định xuyên văn hóa \cite{lee2025globalmood} song chưa có chuẩn màu--cảm xúc bản địa cho Việt Nam; (iii) màu ấm và bão hòa (nâu, hồng) được đọc là ``năng lượng cao'' --- điều này \textit{trung thành} với quan hệ màu--âm nhạc của Whiteford dù có thể nghịch trực giác.

\subsection{Đánh giá chức năng gợi ý bài tương tự}
\label{subsection:eval-similar}
Chức năng này được đánh giá trên các độ đo nội tại (80 bài hạt giống) và trên độ lợi NDCG \cite{steck2018calibrated} đối chiếu với các danh sách phát biên tập, theo tinh thần dùng baseline mạnh trong đánh giá hệ gợi ý \cite{dacrema2019progress}. Độ nhất quán tâm trạng (MoodCoherence) đo mức các bài kết quả cụm chặt trong không gian V-A, định nghĩa theo Phương trình~\eqref{eq:moodcoh}; độ nhất quán nhịp độ (TempoCoherence) định nghĩa tương tự trên nhịp độ.
\begin{equation}\label{eq:moodcoh}
    \text{MoodCoherence} = 1 - \frac{2}{k(k-1)}\sum_{i<j} \big\lVert \mathbf{p}_i - \mathbf{p}_j \big\rVert_2
\end{equation}
Về độ đo nội tại (80 bài hạt giống, $top\text{-}k=10$), MoodCoherence đạt $\approx 0.94$ và TempoCoherence đạt $\approx 0.84$ --- các bài kết quả cụm chặt cả về cảm xúc lẫn nhịp độ.

\textbf{Ground truth tương đồng âm nhạc (độ đo chính).} Để đánh giá đúng mức ``nghe giống nhau'', đồ án xây dựng một tập ground truth \textit{tương đồng âm nhạc} gồm 52 bài hạt giống, mỗi cặp được chấm theo thang liên quan $0$--$3$. Một lo ngại với ground truth do mô hình sinh là nó có thể là tạo tác của \textit{một} bộ chấm duy nhất; vì vậy toàn bộ 1.048 cặp được chấm lại độc lập bằng \textit{hai mô hình của hai nhà cung cấp khác nhau} (qwen3:8b cục bộ và GPT-4o-mini). Hai bộ chấm đồng thuận gần như tuyệt đối (Bảng~\ref{table:multijudge}): Spearman $\rho = 0.946$, đồng thuận nhị phân $0.99$ và Cohen $\kappa = 0.958$ --- xác nhận ground truth đáng tin chứ không phải hiện tượng riêng của một bộ chấm. Trên tập này, hệ thống đạt NDCG@10 kiểm-chứng-chéo $\approx 0.73$ (chi tiết khảo sát backbone ở mục~\ref{subsection:eval-bakeoff}).

\textbf{NDCG trên danh sách phát biên tập (độ đo phụ, đọc kèm baseline).} Trên 872 truy vấn từ danh sách phát biên tập, đo bằng độ lợi tích lũy chiết khấu chuẩn hóa NDCG \cite{jarvelin2002ndcg} (Phương trình~\eqref{eq:ndcg}), hệ sản phẩm đạt NDCG@10 $\approx 0.066$ (tái lập bằng \texttt{tools/eval\_editorial\_ndcg.py}). Tuy nhiên độ đo này \textit{không phù hợp} làm thước đo chính cho gợi ý bài tương tự: một baseline \textit{theo độ phổ biến} (chọn 10 bài theo tần suất nghệ sĩ) đạt $0.091$ --- \textit{cao hơn} hệ sản phẩm --- trong khi baseline ngẫu nhiên chỉ đạt $0.036$ (Bảng~\ref{table:editorial}). Đây không phải do hệ kém: danh sách phát biên tập \textit{thưởng} cho độ phổ biến và đồng-xuất-hiện, đồng thời \textit{phạt} chính cơ chế đa dạng nghệ sĩ (MMR) và lọc bản cover mà hệ cố ý áp dụng --- đúng như cảnh báo về việc dùng độ đo đại diện yếu trong đánh giá hệ gợi ý \cite{dacrema2019progress}. Do đó đồ án dùng ground truth tương đồng âm nhạc cùng các độ đo nội tại làm thước đo chính, còn NDCG biên tập chỉ để tham khảo xu hướng.
\begin{equation}\label{eq:ndcg}
    \text{NDCG@}k = \frac{\text{DCG@}k}{\text{IDCG@}k}, \qquad
    \text{DCG@}k = \sum_{i=1}^{k} \frac{rel_i}{\log_2(i+1)}
\end{equation}

\begin{table}[h!]
\centering
\caption{Đồng thuận hai bộ chấm độc lập trên ground truth tương đồng âm nhạc (1.048 cặp)}
\label{table:multijudge}
\begin{tabular}{p{8.0cm} c}
\toprule
\textbf{Chỉ số đồng thuận qwen3:8b $\leftrightarrow$ GPT-4o-mini} & \textbf{Giá trị} \\
\midrule
Spearman $\rho$ (thang liên quan $0$--$3$) & 0.946 \\
Đồng thuận nhị phân (liên quan $\geq 2$) & 0.99 \\
Cohen $\kappa$ & 0.958 \\
\bottomrule
\end{tabular}
\end{table}

\begin{table}[h!]
\centering
\caption{NDCG@10 trên danh sách phát biên tập kèm baseline (đọc kèm cảnh báo Dacrema)}
\label{table:editorial}
\begin{tabular}{p{6.0cm} c c}
\toprule
\textbf{Phương pháp} & \textbf{NDCG@10} & \textbf{$\times$ ngẫu nhiên} \\
\midrule
Theo độ phổ biến (baseline) & 0.091 & 2.5$\times$ \\
Sản phẩm (Brightify) & 0.066 & 1.8$\times$ \\
Ngẫu nhiên (baseline) & 0.036 & 1.0$\times$ \\
\bottomrule
\end{tabular}
\end{table}

Khả năng chuyển miền của các đầu dò cảm xúc được kiểm trên tập PMEmo \cite{zhang2018pmemo} (Hình~\ref{fig:pmemo}), đạt tương quan $0.69$ cho Valence và $0.65$ cho Arousal, cho thấy đầu dò nắm bắt cảm xúc thật chứ không phải hiện tượng riêng của tập huấn luyện; một kiểm chứng arousal độc lập bằng đặc trưng âm thanh khác (CLAP \cite{wu2023clap}) cho tương quan $0.46$.

\begin{figure}[h!]
\centering
\includegraphics[width=0.5\textwidth]{Hinhve/fig_pmemo.png}
\caption{Chuyển miền các đầu dò cảm xúc sang tập PMEmo độc lập (tương quan Spearman, kèm khoảng tin cậy)}
\label{fig:pmemo}
\end{figure}

\subsection{Khảo sát so sánh mô hình}
\label{subsection:eval-bakeoff}
Để bảo đảm mỗi thành phần học máy là lựa chọn tốt nhất, đồ án tiến hành một khảo sát so sánh trên \textbf{năm vị trí có thể thay mô hình}, với tổng cộng hơn mười mô hình tiên tiến được tải về và chạy thực tế; mỗi ứng viên được tinh chỉnh siêu tham số riêng trước khi so và được phán quyết bằng độ đo cuối theo phương pháp ở mục~\ref{subsection:eval-method}. Bảng~\ref{table:bakeoff} tóm tắt toàn bộ khảo sát, còn Bảng~\ref{table:backbone} trình bày sâu một head-to-head công bằng tiêu biểu (MuQ so với MERT). Điểm đáng chú ý xuyên suốt là phần lớn các mô hình ``mạnh trên giấy'' theo bộ chuẩn quốc tế lại \textit{không} vượt được độ đo cuối của hệ thống: MuQ-MuLan tối ưu cho việc ghép văn bản--nhạc nên cho sai số màu và độ nhất quán kém hơn MuQ; multilingual-e5-large \cite{wang2024e5} thắng các mô hình embedding lời còn lại (kể cả mô hình tiếng Việt tiền nhiệm) về độ nhất quán nội tại ở mọi trọng số; PhoBERT-large chỉ ngang ViSoBERT (ViSoBERT gọn hơn và hợp văn phong lời nhạc hơn); mDeBERTa-v3 kém XLM-R ở vai trò đầu dò liên ngữ; và với trục Arousal của màu, mô hình màu--nhạc của Whiteford được chọn thay cho Valdez--Mehrabian và bản hồi quy theo chuẩn ICEAS vì hai bản sau khớp arousal \textit{màu khi nhìn} (coi màu tối là kích thích) --- cấu trúc sai cho một hệ gợi ý nhạc. Kết quả này củng cố nguyên tắc đánh giá bằng độ đo cuối thay vì thứ hạng trên bộ chuẩn. Ngoài các thành phần phục vụ, hệ thống còn dùng CLAP làm một kiểm chứng arousal độc lập (mục~\ref{subsection:eval-similar}) và đã thử nhưng \textit{loại} bộ đặc trưng Essentia do đặc trưng suy biến ở tần số lấy mẫu 44.1\,kHz. Chi tiết phân tích khảo sát này được trình bày ở Chương~\ref{chapter:SolutionAndContribution}.

\textbf{Lựa chọn trục biểu diễn âm thanh (MuQ so với MERT).} Việc chọn MuQ thay cho MERT được kiểm bằng một so sánh \textit{công bằng}: chuyển backbone thật và \textit{tối ưu lại trọng số tổ hợp riêng cho từng backbone} bằng kiểm chứng chéo 5-fold trên ground truth tương đồng âm nhạc (Bảng~\ref{table:backbone}). Trên các độ đo trực tiếp, hai backbone gần như \textit{ngang nhau} (NDCG kiểm-chứng-chéo $0.735$ so với $0.734$; NDCG chỉ-âm-thanh $0.748$ so với $0.742$ --- chênh lệch không có ý nghĩa thống kê), và sai số màu TE coi như hòa. Khác biệt có ý nghĩa nằm ở \textit{phép chuyển giao sang một bộ chấm độc lập} (GPT-4o-mini): ở đó MuQ vượt MERT $0.806$ so với $0.771$ (khoảng tin cậy bootstrap cặp $[+0.014, +0.056]$, có ý nghĩa thống kê). Cùng với việc MuQ thắng ở đầu dò Arousal trên DEAM ($0.69$ so với $0.647$, mục~\ref{subsection:4.2.1}), kết quả cho thấy MuQ \textit{ngang hoặc hơn} MERT ở mọi mặt và việc chọn MuQ là hợp lý. Lưu ý phương pháp: trọng số tổ hợp tối ưu theo kiểm chứng chéo hội tụ về gần như chỉ-âm-thanh; trọng số sản phẩm $0.76/0.16/0.08$ đánh đổi rất ít NDCG để tăng độ nhất quán cảm xúc.

\begin{table}[h!]
\centering
\caption{So sánh công bằng MuQ với MERT (mỗi backbone tối ưu trọng số riêng)}
\label{table:backbone}
\begin{tabular}{p{7.6cm} c c}
\toprule
\textbf{Phép đo} & \textbf{MuQ} & \textbf{MERT} \\
\midrule
NDCG@10 kiểm-chứng-chéo (GT tương đồng âm nhạc) & 0.735 & 0.734 \\
NDCG@10 chỉ-âm-thanh (cô lập backbone) & 0.748 & 0.742 \\
NDCG@10 chuyển sang bộ chấm độc lập (GPT) & \textbf{0.806} & 0.771 \\
Tương quan Arousal kiểm-chứng-chéo trên DEAM & \textbf{0.69} & 0.647 \\
\bottomrule
\end{tabular}
\end{table}

\begin{table}[h!]
\centering
\caption{Khảo sát so sánh mô hình trên năm vị trí (mô hình được chọn in đậm)}
\label{table:bakeoff}
\begin{tabular}{p{2.2cm} p{2.5cm} p{3.3cm} p{4.5cm}}
\toprule
\textbf{Vị trí} & \textbf{Được chọn} & \textbf{Ứng viên đã so sánh} & \textbf{Phán quyết trên độ đo cuối} \\
\midrule
Biểu diễn âm thanh & \textbf{MuQ} & MERT, MuQ-MuLan & MuQ $\geq$ MERT (Bảng~\ref{table:backbone}); MuLan kém sai số màu và nhất quán vì tối ưu ghép văn bản--nhạc \\
Embedding lời & \textbf{multilingual\hspace{0pt}-e5-large} & VN-SBERT (tiền nhiệm), BGE-M3, PhoBERT-large, mDeBERTa & Thắng độ nhất quán nội tại ở 200 hạt giống, mọi trọng số \\
Cảm xúc lời (VN) & \textbf{ViSoBERT} & PhoBERT-large, wonrax-sentiment, ViDeBERTa & PhoBERT-large chỉ ngang (sau khi bật tách từ); ViSoBERT gọn và hợp văn phong lời hơn \\
Cảm xúc liên ngữ & \textbf{XLM-R} & mDeBERTa-v3 & mDeBERTa + gộp trung bình vẫn kém CCC \\
Màu $\rightarrow$ Arousal & \textbf{Whiteford 2018} & Valdez--Mehrabian, hồi quy ICEAS & Whiteford đúng cấu trúc màu--nhạc (kiểm bằng BPM thật); hai bản kia khớp arousal màu khi nhìn nên màu tối hóa ``năng lượng cao'' sai \\
\bottomrule
\end{tabular}
\end{table}

\subsection{Phân tích độ nhạy của trọng số}
\label{subsection:eval-sensitivity}
Một phần trọng số trong hệ thống không được khớp trực tiếp từ dữ liệu người gán mà được \textit{đặt theo nguyên tắc} (trọng số nhịp độ trong Arousal, băng thông $\sigma$ của RBF màu) hoặc \textit{đánh đổi có chủ đích} giữa hai mục tiêu (trọng số V-A trong tổ hợp gợi ý bài tương tự). Để chứng minh các giá trị đang phục vụ không phải lựa chọn tùy tiện, đồ án thực hiện một phân tích độ nhạy: quét quanh mỗi giá trị và quan sát độ đo cuối thay đổi ra sao. Phân tích này \textit{không} thay đổi gì trong hệ đang chạy --- nó chỉ đo trên bộ nhớ và có thể tái lập bằng \texttt{tools/tune\_muq\_arousal.py} và \texttt{tools/sensitivity\_analysis.py}.

\textbf{Trọng số nhịp độ (Arousal).} Đây là một \textit{đánh đổi}, không phải điểm phẳng: tăng trọng số nhịp độ giúp Arousal bám nhịp độ thật tốt hơn nhưng làm giảm độ chính xác kiểm-chứng-chéo trên DEAM (Hình~\ref{fig:sens-tempo}). Giá trị phục vụ $0.35$ nằm trong vùng hợp lệ: nó đạt $\rho(\text{Arousal},\text{BPM}) = 0.47$ (vượt xa mục tiêu $0.20$) trong khi vẫn giữ DEAM-CV $= 0.69$, cao hơn cổng không-thụt-lùi $0.647$. Tăng lên $0.45$ sẽ phá cổng này (DEAM-CV tụt còn $0.61$), nên cửa sổ hợp lệ là khoảng $[0.25,\,0.40]$ và $0.35$ được chọn để bám nhịp độ mạnh nhất trong cửa sổ đó.

\begin{figure}[h!]
\centering
\includegraphics[width=0.62\textwidth]{Hinhve/fig_sens_tempo.png}
\caption{Độ nhạy của trọng số nhịp độ: vùng xanh là cửa sổ hợp lệ (DEAM-CV $\geq 0.647$ và $\rho(A,\text{BPM}) \geq 0.20$); giá trị phục vụ $0.35$ nằm trong vùng này}
\label{fig:sens-tempo}
\end{figure}

\textbf{Băng thông $\sigma$ của RBF màu.} Ngược lại, kết quả gợi ý theo màu \textit{rất bền} với lựa chọn $\sigma$: khi quét $\sigma_V$ trong $[0.12,\,0.28]$ và $\sigma_A$ trong $[0.08,\,0.20]$, sai số nhắm đích (TE) chỉ dao động chưa tới $0.003$ (Hình~\ref{fig:sens-sigma}). Điều này cho thấy giá trị phục vụ ($\sigma_V = 0.20$, $\sigma_A = 0.14$) --- được đặt theo tỉ lệ độ tin cậy của hai trục (arousal đáng tin hơn nên $\sigma$ hẹp hơn) --- không phải tham số mong manh; kết luận không phụ thuộc vào con số chính xác.

\begin{figure}[h!]
\centering
\includegraphics[width=0.62\textwidth]{Hinhve/fig_sens_sigma.png}
\caption{Độ nhạy của băng thông $\sigma$ (RBF màu): TE gần như phẳng trên toàn lưới quét, xác nhận hệ bền với lựa chọn $\sigma$}
\label{fig:sens-sigma}
\end{figure}

\textbf{Trọng số V-A trong gợi ý bài tương tự.} Đây là một đánh đổi được công bố thẳng. Khi tăng trọng số V-A từ $0$, NDCG@10 (trên ground truth tương đồng âm nhạc) đạt cao nhất ở mức gần \textit{chỉ-âm-thanh} rồi giảm dần, trong khi độ nhất quán cảm xúc (MoodCoherence) tăng mạnh (Hình~\ref{fig:sens-fusion}, Bảng~\ref{table:sens-fusion}). Giá trị phục vụ $0.16$ đúng là điểm \textit{khuỷu}: đi từ $0$ lên $0.16$, MoodCoherence nhảy từ $0.83$ lên $0.95$ ($+0.12$) trong khi NDCG chỉ giảm $\approx 0.003$; vượt quá $0.16$ thì NDCG tiếp tục giảm còn MoodCoherence gần như bão hòa. Nói cách khác, $0.16$ thu gần như toàn bộ phần lợi về nhất quán cảm xúc với chi phí NDCG rất nhỏ. (Lưu ý: giá trị NDCG tuyệt đối ở đây thấp vì được tính bằng cách xếp hạng \textit{toàn bộ kho nhạc} rồi đối chiếu với tập đã chấm nhỏ, khác với NDCG kiểm-chứng-chéo trên pool ở mục~\ref{subsection:eval-bakeoff}; chỉ \textit{xu hướng tương đối} giữa các trọng số là điều cần đọc.)

\begin{table}[h!]
\centering
\caption{Độ nhạy của trọng số V-A trong gợi ý bài tương tự (52 hạt giống ground truth tương đồng âm nhạc)}
\label{table:sens-fusion}
\small
\begin{tabular}{c c c c}
\toprule
\textbf{Trọng số V-A} & \textbf{Trọng số âm thanh} & \textbf{NDCG@10 (xu hướng)} & \textbf{MoodCoherence} \\
\midrule
0.00 & 0.92 & 0.0414 & 0.833 \\
0.08 & 0.84 & 0.0362 & 0.924 \\
\textbf{0.16 (phục vụ)} & \textbf{0.76} & \textbf{0.0386} & \textbf{0.948} \\
0.24 & 0.68 & 0.0335 & 0.957 \\
0.32 & 0.60 & 0.0299 & 0.964 \\
\bottomrule
\end{tabular}
\end{table}

\begin{figure}[h!]
\centering
\includegraphics[width=0.66\textwidth]{Hinhve/fig_sens_fusion.png}
\caption{Đánh đổi NDCG $\leftrightarrow$ MoodCoherence theo trọng số V-A: giá trị phục vụ $0.16$ là điểm khuỷu --- gần như toàn bộ lợi ích nhất quán cảm xúc đạt được với chi phí NDCG nhỏ}
\label{fig:sens-fusion}
\end{figure}

Tổng hợp lại, phân tích độ nhạy cho thấy: ở những trục đặt theo nguyên tắc ($\sigma$), kết quả bền với lựa chọn chính xác; còn ở những trục mang tính đánh đổi (nhịp độ, V-A), giá trị phục vụ nằm đúng vùng/điểm cân bằng đã nêu lý do. Không có trọng số nào được chọn cảm tính mà không kiểm chứng vùng lân cận.

\subsection{Bàn luận và mối đe dọa tính hợp lệ}
\label{subsection:eval-discussion}
Tổng hợp lại, các kết quả trên cho thấy hệ thống đạt mục tiêu kép: khớp cảm xúc (TE thấp, vượt mọi baseline thực) và nghe đồng nhất (độ nhất quán tăng mạnh). Tuy nhiên, để diễn giải đúng mức ý nghĩa của chúng, cần nêu rõ các mối đe dọa tính hợp lệ theo ba nhóm thường dùng trong đánh giá hệ thống.

\textit{Hợp lệ về cấu trúc} (construct validity) --- liệu độ đo có đo đúng thứ cần đo. Sai số nhắm đích TE đo trên chính nhãn V-A được dùng để truy hồi, nên về bản chất là \textit{tự tham chiếu}: nó kiểm tra hệ thống có truy hồi đúng theo định nghĩa của chính nó hay không, chứ chưa khẳng định nhãn V-A là ``đúng'' theo cảm nhận người. Đây là lý do đồ án bổ sung các kiểm chứng \textit{độc lập với nhãn}: tương quan với nhịp độ thật (cho trục Arousal của màu), khớp chuẩn ICEAS (cho trục Valence của màu), và chuyển miền sang PMEmo (cho các đầu dò cảm xúc). Sự nhất quán giữa các kiểm chứng độc lập này làm tăng đáng kể độ tin cậy rằng nhãn V-A nắm bắt cảm xúc thật.

\textit{Nguyên tắc dùng dữ liệu nước ngoài.} Một câu hỏi tự nhiên là vì sao có thể dùng các tập do con người gán nhãn ở nước ngoài (DEAM, EmoBank, PMEmo, chuẩn màu ICEAS) cho một hệ thống nhạc Việt. Nguyên tắc xuyên suốt là dữ liệu nước ngoài chỉ được dùng ở những trục mà tín hiệu \emph{không phụ thuộc ngôn ngữ}. Arousal được suy từ năng lượng âm thanh (nhịp độ, độ to, đặc trưng MuQ) --- một bài nhanh và to mang năng lượng cao bất kể hát bằng tiếng nào --- nên đầu dò huấn luyện trên DEAM (nhạc phương Tây) chuyển sang nhạc Việt là hợp lệ; tương tự, ánh xạ màu dựa trên tri giác màu của chuẩn ICEAS. Ngược lại, Valence \emph{phụ thuộc nghĩa của lời} nên \emph{không} được lấy chủ yếu từ dữ liệu nước ngoài: trong tổ hợp EWE, các tín hiệu lời chiếm $\approx 91\%$ trọng số, trong đó $\approx 57\%$ đến từ tài nguyên \emph{tiếng Việt bản địa} (từ điển NRC-VAD-VN cùng VnEmoLex, và đầu dò ViSoBERT trên dữ liệu cảm xúc tiếng Việt) và $\approx 34\%$ từ EmoBank tiếng Anh được \emph{chuyển ngữ} qua mô hình đa ngữ XLM-R; đầu dò âm thanh chỉ đóng góp $\approx 9\%$. Cách phân vai theo phương thức (âm thanh/tri giác dùng nguồn quốc tế, ngữ nghĩa lời dùng nguồn bản địa) chính là điều được kiểm chứng bằng phép đo nêu ở mục hợp lệ ngoại tại: nếu cố suy Valence từ âm thanh nước ngoài thì kết quả rất yếu ($R^2 \approx 0.06$).

\textit{Vì sao không tự gán nhãn cảm xúc cho kho nhạc Việt.} Việc không xây một bộ nhãn cảm xúc do chính tác giả gán là một lựa chọn có chủ đích về phương pháp, không phải do thiếu dữ liệu. Thứ nhất, một bộ nhãn do \emph{một người gán duy nhất} không có độ tin cậy liên-người-gán (inter-rater reliability) nên về mặt thống kê không cấu thành một ground truth đáng tin --- nó chỉ phản ánh phán đoán của một cá nhân. Thứ hai, vì người gán cũng chính là người xây dựng hệ thống, nhãn tự gán có nguy cơ \emph{thiên lệch xác nhận} (confirmation bias) khiến việc ``kiểm chứng'' trở thành lập luận vòng tròn. Thay vì tạo ra một chuẩn so yếu và dễ thiên lệch như vậy, đồ án kiểm chứng nhãn bằng sự \emph{hội tụ với nhiều tập do con người gán nhãn đã công bố và bình duyệt} (DEAM, PMEmo cho âm thanh; ICEAS cho màu --- mỗi tập là phán đoán tổng hợp của hàng chục đến hàng trăm người gán độc lập) cùng một panel nhiều bộ chấm (mục~\ref{subsection:eval-similar}); riêng Valence còn được đối chiếu với một bộ tham chiếu GPT-4o-mini đọc trực tiếp lời tiếng Việt ($\rho = 0.63$, mục~\ref{subsection:4.2.1}). Sự nhất quán giữa các nguồn độc lập này là bằng chứng khách quan hơn so với một tập tự gán đơn lẻ. Việc xây dựng một bộ dữ liệu cảm xúc nhạc Việt có nhiều người gán và báo cáo độ tin liên-người-gán được để lại như một hướng phát triển (Chương~\ref{chapter:conclusion}).

\textit{Hợp lệ nội tại} (internal validity) --- liệu kết luận so sánh có vững. Việc so với năm baseline (trong đó có hai đường loại bỏ từng trục) kèm kiểm định Wilcoxon và hiệu chỉnh FDR bảo đảm cải thiện không phải do ngẫu nhiên; khoảng tin cậy bootstrap cho thấy độ ổn định của ước lượng. Riêng đường ``lân cận V-A tham lam'' là một trần lý thuyết (oracle) chứ không phải đối thủ công bằng, nên việc hệ thống bám sát nó (chứ không vượt) là kết quả mong đợi, cho thấy phần đa dạng hóa chỉ đánh đổi rất ít độ chính xác nhắm đích. Riêng với gợi ý bài tương tự, ground truth tương đồng âm nhạc còn được xác nhận bằng hai bộ chấm độc lập của hai nhà cung cấp ($\kappa = 0.958$, mục~\ref{subsection:eval-similar}), nên kết luận so sánh không phải là tạo tác của một bộ chấm duy nhất; hạn chế còn lại là cỡ mẫu hạt giống ($n = 52$).

\textit{Hợp lệ ngoại tại} (external validity) --- liệu kết quả có tổng quát hóa. Đây là hạn chế lớn nhất: chưa có nghiên cứu nghe của người dùng thật (theo Dacrema và cộng sự, tương quan giữa độ đo ngoại tuyến và mức hài lòng thực tế chỉ vào khoảng $r \approx 0.28$), và ánh xạ màu--cảm xúc là ngoại suy từ dữ liệu quốc tế (Việt Nam không nằm trong khảo sát ICEAS gốc). Một hạn chế cụ thể của trục Valence là \textit{chuyển giao yếu qua biên giới văn hóa}: trong phạm vi nhạc phương Tây, phép chuyển miền DEAM$\to$PMEmo khá tốt ($\rho = 0.69$ cho Valence), nhưng khi áp đầu dò âm thanh huấn luyện trên DEAM sang nhạc Việt thì Valence rất yếu ($R^2 \approx 0.06$). Đây chính là lý do Valence được lấy chủ yếu từ các tín hiệu lời (phần lớn là tài nguyên tiếng Việt bản địa, phần còn lại là EmoBank tiếng Anh chuyển ngữ qua XLM-R; chi tiết phân vai và trọng số nêu ở phần hợp lệ cấu trúc trên) thay vì từ âm thanh nước ngoài --- một lựa chọn đúng nhưng phụ thuộc chất lượng các tài nguyên đó và độ tốt của phép chuyển ngữ liên ngữ. Ngoài ra, do tương quan dương giữa Valence và Arousal trong kho nhạc, các góc phần tư lệch trục (``bình yên'': V cao--A thấp; ``căng thẳng'': V thấp--A cao) thưa dữ liệu, khiến màu ấm phải nhắm vào phần đuôi Arousal cao mỏng ($\approx 22\%$ kho nhạc) --- đây là thuộc tính của dữ liệu. Chỉ số NDCG@10 biên tập tuyệt đối thấp ($\approx 0.066$) cũng cần đọc thận trọng: nó là đại diện gần đúng cho mức liên quan và còn bị một baseline theo độ phổ biến vượt qua (mục~\ref{subsection:eval-similar}), nên chỉ phù hợp để tham khảo xu hướng chứ không khẳng định chất lượng tuyệt đối. Những hạn chế này được nêu lại và đề xuất hướng khắc phục ở Chương~\ref{chapter:conclusion}.

\section{Thời gian phản hồi và triển khai}
\label{section:deploy}
Hệ thống được triển khai thử nghiệm với pha ngoại tuyến chạy trước để sinh các đặc trưng và ma trận biểu diễn cho toàn bộ kho nhạc, sau đó động cơ gợi ý nạp các đặc trưng này vào bộ nhớ và phục vụ truy vấn qua giao diện lập trình ứng dụng cho ứng dụng trang đơn.

\subsection{Thời gian phản hồi đo được}
\label{subsec:latency}
Vì mọi tính toán nặng đã thực hiện ở pha ngoại tuyến và không có mô hình nào được suy luận lúc phục vụ (nhãn cảm xúc đã đóng băng thành tra cứu), mỗi truy vấn chỉ gồm vài phép nhân ma trận trên ma trận $5138 \times 1024$. Đo trực tiếp trong tiến trình trên một máy Apple Silicon bằng \texttt{tools/bench\_latency.py} (1500 lần lặp, chưa tính thời gian mạng và tuần tự hóa JSON), kết quả như Bảng~\ref{table:latency} và Hình~\ref{fig:latency}: nạp động cơ lõi một lần chỉ mất vài giây (đo được khoảng $1$--$6$ giây tùy bộ nhớ đệm; khởi động đầy đủ kèm bộ mã hóa tìm kiếm tùy chọn sẽ lâu hơn), còn mỗi truy vấn lúc phục vụ ở mức \textbf{vài mili-giây} (gợi ý theo một màu $\approx 3.1$ ms, gợi ý bài tương tự $\approx 4.1$ ms); hành trình hai màu nặng hơn ($\approx 70$ ms) do bước sắp xếp chuỗi cảm xúc. Đây là độ trễ tính toán lõi; độ trễ đầu cuối nhìn từ trình duyệt sẽ cao hơn do mạng và phát âm thanh.

\begin{table}[h!]
\centering
\caption{Thời gian phản hồi tính toán lõi (đo trong tiến trình, không gồm mạng)}
\label{table:latency}
\begin{tabular}{p{6.6cm} c c}
\toprule
\textbf{Thao tác} & \textbf{Trung vị} & \textbf{p95} \\
\midrule
Nạp động cơ lõi (một lần, lúc khởi động) & $\approx 1$--$6$ s & --- \\
Gợi ý theo một màu (UC02) & $3.1$ ms & $3.6$ ms \\
Gợi ý bài tương tự (UC01) & $4.1$ ms & $4.5$ ms \\
Hành trình hai màu & $70$ ms & $72$ ms \\
\bottomrule
\end{tabular}
\end{table}

\begin{figure}[h!]
\centering
\includegraphics[width=0.68\textwidth]{Hinhve/fig_latency.png}
\caption{Thời gian phản hồi tính toán lõi theo thao tác (thang log); truy vấn đơn ở mức vài mili-giây}
\label{fig:latency}
\end{figure}

\section{Kiểm thử chức năng hệ thống}
\label{section:functest}
Ngoài đánh giá mô hình, hệ thống còn được kiểm thử như một sản phẩm phần mềm. Bảng~\ref{table:testcase} tóm tắt các ca kiểm thử chính. Một phần được tự động hóa bằng \texttt{pytest} (\texttt{test/test\_color\_reco.py} kiểm tính hợp lệ cấu hình, ánh xạ 12 màu về đúng góc phần tư cảm xúc, đa dạng nghệ sĩ, tính đơn điệu của hành trình, độ trải Arousal; \texttt{test/test\_endless\_radio.py} kiểm hợp đồng loại trừ để chế độ ``đài'' không lặp), một phần được kiểm theo hợp đồng API và thủ công. Tổng cộng 37 ca kiểm thử tự động trong hai tệp cốt lõi này hiện đều đạt.

\begin{table}[h!]
\centering
\caption{Các ca kiểm thử chức năng hệ thống}
\label{table:testcase}
\begin{tabular}{p{1.0cm} p{4.4cm} p{5.2cm} p{2.6cm}}
\toprule
\textbf{Mã} & \textbf{Mục tiêu} & \textbf{Kết quả mong đợi} & \textbf{Trạng thái} \\
\midrule
TC01 & Chọn màu và nhận danh sách bài hát & Trả về $\geq 5$ bài, V-A hợp lệ, đa dạng nghệ sĩ & Đạt (tự động) \\
TC02 & Chọn bài và nhận bài tương tự & Trả về danh sách không chứa bài nguồn, đã loại cover & Đạt (tự động) \\
TC03 & Xử lý mã màu không hợp lệ & Từ chối với lỗi xác thực đầu vào, không lỗi máy chủ & Đạt (hợp đồng API) \\
TC04 & Xử lý bài không có lời & Trường lời rỗng, không gây lỗi & Đạt (hợp đồng API) \\
TC05 & Xử lý bài không có đặc trưng & Trường tương ứng rỗng; lùi về tín hiệu còn lại & Đạt (hợp đồng API) \\
TC06 & Kiểm tra thời gian phản hồi & Truy vấn đơn ở mức vài mili-giây & Đạt (mục~\ref{subsec:latency}) \\
TC07 & Kiểm tra giao diện trên trình duyệt & Hai giao diện dựng và hoạt động & Đạt một phần (smoke; đa trình duyệt: hướng hoàn thiện) \\
\bottomrule
\end{tabular}
\end{table}

Chương này đã trình bày kiến trúc, thiết kế chi tiết, dữ liệu và kho nhạc, giao diện lập trình ứng dụng và giao diện người dùng, kết quả đánh giá có thể kiểm chứng, thời gian phản hồi đo được và phần kiểm thử chức năng. Những giải pháp và đóng góp nổi bật rút ra từ quá trình này được phân tích trong chương tiếp theo.

\end{document}
